\chapter{Parte 2: sistema anticolisiones}

La segunda parte modifica el comportamiento anterior para que el robot esquive obstaculos mientras mantiene atraccion hacia el objetivo. Se añade una fuerza repulsiva asociada a sensores de proximidad:

\[
\vec{F}_{rep} =
\begin{cases}
k_{rep}\left(\frac{1}{d}-\frac{1}{d_0}\right)\frac{1}{d^2}\hat{u}_{obs}, & d < d_0\\
0, & d \ge d_0
\end{cases}
\]

donde \(d\) es la distancia al obstaculo, \(d_0\) el umbral de influencia y \(\hat{u}_{obs}\) la direccion de alejamiento.

\section{Combinacion de fuerzas}

La fuerza total se calcula como:

\[
\vec{F}_{total} = \vec{F}_{att} + \sum_i \vec{F}_{rep,i}
\]

El vector resultante determina el rumbo deseado. Si un obstaculo aparece frontalmente, la componente repulsiva modifica el angulo de giro y reduce la velocidad. Para evitar movimientos inestables se aplican saturaciones y un filtro simple sobre la consigna.

\section{Casos de prueba}

\begin{itemize}
  \item Obstaculo aislado entre robot y objetivo.
  \item Pasillo estrecho con margen suficiente.
  \item Objetivo movil representado por Bill.
  \item Obstaculo inesperado durante la aproximacion.
  \item Escenario sin obstaculos para comprobar que no se degrada la Parte 1.
\end{itemize}
